% ============================================================
% Capitolo 10 — HALT Triggers
% ============================================================

\chapter{HALT Triggers: il freno di emergenza}
\label{ch:halt-triggers}

Ci sono zone in ogni progetto che non dovresti mai modificare di impulso: schema del database, codice di autenticazione, credenziali, CI/CD. Gli HALT trigger sono il meccanismo con cui AI-DLC presidia queste zone.

\section{Il file \texttt{halt-triggers.yaml}}
\label{sec:10-file}

Vive in \texttt{.ai-dlc/halt-triggers.yaml} ed è già lì dopo \texttt{init}. I sei trigger predefiniti:

\begin{lstlisting}[language=yaml,caption={Trigger predefiniti in halt-triggers.yaml}]
triggers:
  - id: schema
    patterns: ["**/migrations/**", "**/prisma/schema.prisma",
               "**/*.sql", "**/alembic/**"]
    risk: HIGH

  - id: auth
    patterns: ["**/auth/**", "**/middleware/auth*",
               "**/oauth/**"]
    risk: HIGH+

  - id: secrets
    patterns: ["**/.env*", "**/secrets/**", "**/*.pem"]
    risk: HIGH+

  - id: infra
    patterns: ["**/terraform/**", "**/k8s/**",
               "**/Dockerfile*", "**/*.tf"]
    risk: HIGH

  - id: cicd
    patterns: [".github/workflows/**", "**/.gitlab-ci.yml"]
    risk: HIGH

  - id: framework
    patterns: [".ai-dlc/modules/**", "AGENTS.md", "CLAUDE.md"]
    risk: HIGH
\end{lstlisting}

\section{Override di progetto}
\label{sec:10-override}

Le zone sensibili specifiche del tuo progetto vanno in \texttt{.ai-dlc/project/halt-triggers.yaml}:

\begin{lstlisting}[language=yaml,caption={Override di progetto per TaskFlow API}]
version: 1

triggers:
  - id: billing
    patterns: ["**/billing/**", "**/payments/**"]
    reason: Logica di pagamento - impatto finanziario
    risk: HIGH+
\end{lstlisting}

Il file di progetto \textit{aggiunge} trigger a quelli predefiniti. Per disattivare un trigger predefinito temporaneamente:

\begin{lstlisting}[language=yaml,caption={Disattivazione temporanea del trigger schema}]
version: 1
disable:
  - schema  # disattivato durante la fase di design intensivo DB
\end{lstlisting}

\begin{warnbox}
Un trigger disattivato è una zona senza rete di sicurezza. Documenta in \texttt{PROGRESS.md} perché l'hai disattivato e quando pianifichi di riattivarlo.
\end{warnbox}

\section{I trigger non si disattivano mai con i Modes}
\label{sec:10-modes}

Anche in Mode FAST --- il mode con meno cerimonia --- un path che corrisponde a un trigger viene bloccato e richiede conferma. Il rischio di un percorso auth o di uno schema DB non cambia in base a quanto velocemente vuoi lavorare quel giorno.

\section*{\LabelSummary}

\begin{itemize}
  \item I sei trigger predefiniti coprono: schema DB, auth, secrets, infra, CI/CD, framework AI-DLC.
  \item Gli override di progetto in \texttt{.ai-dlc/project/halt-triggers.yaml} estendono (o disattivano) i trigger predefiniti.
  \item I trigger non si disattivano mai con i Modes --- sono il livello di sicurezza costante.
  \item Un falso positivo (trigger attivato per path poi non modificato) è accettabile e preferibile al contrario.
\end{itemize}
